\chapter{SigmaFlow: The Minimal Conversion}
\label{ch:sigmaflow}

\section{Design principle and what actually changed}
\label{sec:minimal-change}

The comparison in Chapter~\ref{ch:results} is only interpretable if the
substitution is confined to the generative mechanism. This was treated as a
constraint on the implementation rather than an aspiration, and the resulting
partition is given in Table~\ref{tab:minimal-change}.

\begin{table}[htbp]
\centering
\caption{Partition of the system under the minimal-change constraint.
``Identical'' was verified by file hash over the relevant source trees.}
\label{tab:minimal-change}
\small
\begin{tabular}{@{}p{0.30\textwidth} p{0.16\textwidth} p{0.46\textwidth}@{}}
\toprule
\textbf{Component} & \textbf{Status} & \textbf{Note} \\
\midrule
EquiformerV2 backbone      & identical & all backbone source files hash-identical between arms \\
Featurisation, graph construction & identical & inherited \\
Data pipeline, fragmentation & identical & same conformers, same splits \\
Output block ($\ell=1$ head) & identical & one field $f_i$; Eq.~\eqref{eq:force-field} \\
Newton--Euler reduction     & identical & Eqs.~\eqref{eq:reduce-trans}--\eqref{eq:reduce-rot} \\
Optimiser, LR schedule      & identical & AdamW \citep{loshchilov2019adamw}, same cosine schedule \\
Loss weights                & identical & SigmaDock production values \\
\midrule
Forward process             & \textbf{removed} & no noising process; replaced by a probability path \\
Score interpretation and rescaling & \textbf{removed} & the scalar apparatus converting scores to displacements is deleted, not modified \\
Probability path            & \textbf{new} & Eqs.~\eqref{eq:sf-trans}--\eqref{eq:sf-rot} \\
Training target             & \textbf{new} & conditional velocity in place of score \\
Sampler                     & \textbf{new} & deterministic Euler, Eq.~\eqref{eq:euler}, in place of reverse SDE \\
\midrule
Logarithmic map clamp       & \textbf{corrected} & \S\ref{sec:log-defect} \\
Rotational frame convention & \textbf{corrected} & \S\ref{sec:frame-defect} \\
\bottomrule
\end{tabular}
\end{table}

Two residual confounders should be recorded rather than glossed. The two arms
run in separate software environments with different interpreter versions, and
the SigmaDock arm is invoked with its own rotational-score configuration
flags. Neither touches the backbone, but neither is strictly nothing, and
\S\ref{sec:controlled} returns to them.

\section{A frame defect in the rotational field}
\label{sec:frame-defect}

The convention fixed in \S\ref{sec:trivialisation} requires that the network's
rotational output be a \emph{body-frame} quantity, since the integrator
right-multiplies by its exponential. The Newton--Euler reduction
\eqref{eq:reduce-rot}, however, computes a torque from atom positions
expressed in the global frame, and therefore returns a spatial-frame vector.
In the original implementation this spatial quantity was passed directly to
the integrator, mixing \eqref{eq:body-frame} with \eqref{eq:spatial-frame}.

The correction is the adjoint transport \eqref{eq:adjoint}: the predicted
rotational velocity must be conjugated into the body frame,
\begin{equation}
  \skewm{\Omega}_{\mathrm{body}} \;=\; \Rot_t^\top\,\skewm{w}_{\mathrm{spatial}}\,\Rot_t ,
  \label{eq:frame-fix}
\end{equation}
before it enters $\Exp(\cdot\,\Delta t)$. Two facts make this worth a section
rather than a footnote. First, the error is invisible to any check that does
not vary $\Rot_t$: at $\Rot_t = I$ the two frames agree exactly, so a unit
test at the identity passes. Second, the error is not small --- taking an
Euler step of length $1-t$ under the correct convention reaches $\Rot_1$ to
within $5\times10^{-4}$ degrees, whereas the incorrect left-multiplication
misses by roughly $\Ang{179.5}$.

A terminological hazard contributed to the defect and is worth stating,
because the literature is not consistent. The implementation's own comments
describe the construction as ``right-trivialised'' while implementing
$\dot{\Rot} = \Rot\skewm{\Omega}$, which is the \emph{left}-trivialised or
body convention in the sense of \citet{marsden1999introduction}. The
mathematics was correct; the label was not. \S\ref{sec:diagnostics} reports
the behavioural signature of the correction, and
Appendix~\ref{app:numerics} the numerical verification.

\section{A clamping defect in the logarithmic map}
\label{sec:log-defect}

The inherited implementation of \eqref{eq:so3-log} clamped the argument of
$\arccos$ to $[-0.99,\,0.99]$, a guard intended to avoid a domain error. The
effect is to misreport every angle above roughly $\Ang{172}$ --- close to one
in ten draws from the Haar measure, by the cut-locus remark in
\S\ref{sec:so3}.

For diffusion this is harmless, because the reverse SDE never requires an
exact round trip $\Exp(\Log(\Rot)) = \Rot$. For the geodesic construction
\eqref{eq:sf-rot} it is not, because the interpolant is defined by exactly
that round trip: a misreported $\Log$ places $\Rot_t$ off the intended
geodesic and corrupts the regression target for a non-negligible fraction of
training pairs. The clamp was narrowed to $[-1+10^{-7},\,1-10^{-7}]$.

This is an instance of a general point worth making once: numerical guards
inherited from one generative formulation are not automatically valid in
another, because the two formulations exercise the geometry differently.

\section{Deterministic integration}

Sampling integrates \eqref{eq:euler} from $t=0$ to $t=1$ over a fixed number
of steps. On the translational factor this is ordinary explicit Euler; on the
rotational factor the update is a retraction, meaning the step is taken along
the manifold rather than in an ambient space followed by a projection. Because
$\Exp$ maps into $\SOthree$ exactly, every iterate is a rotation matrix by
construction and no re-orthogonalisation is required --- a small but real
advantage of working in the group rather than with a matrix parameterisation.

Three properties distinguish this from the reverse-time SDE it replaces.
Sampling is \emph{deterministic}, so the pose is a function of the source
draw alone and the diversity of a $K$-sample set is entirely inherited from
$p_0$; the sampler has no temperature or noise-scale to tune. The number of
network evaluations equals the number of steps and is a free parameter at
inference, decoupled from anything fixed at training time. And because the
conditional paths are straight in $\R^3$ and geodesic on $\SOthree$, the
marginal field they induce is expected to be nearly straight, so few Euler
steps should suffice --- a prediction of the construction rather than a hope.
\S\ref{sec:nfe} tests it over a forty-fold range of step counts.

The step count inherited from SigmaDock is $25$, a value tuned there for a
\emph{stochastic} reverse process. There is no reason for that number to
transfer, and \S\ref{sec:nfe} shows that it does not.

\figplaceholder{fig:paths}
{Conditional probability paths on the two factors of $\Mfd$. Left: the
straight-line translational path. Right: the geodesic rotational path, with
the conditional target drawn at three times.}
{Two panels.
\textbf{Left:} $\R^3$ (drawn as $\R^2$). Source sample $x_0$ from a Gaussian
centred at the pocket origin; data point $x_1$; the straight interpolant with
the constant target $x_1-x_0$ drawn as equal arrows at $t=0.2,0.5,0.8$.
\textbf{Right:} $\SOthree$ drawn as a schematic curved surface (a
three-sphere-quotient cartoon is acceptable and standard). $\Rot_0$ sampled
Haar-uniformly, $\Rot_1$ the crystal orientation, the geodesic between them,
and the body-frame target $\vfield_t^{\mathrm{rot}}$ drawn at the same three
times --- \emph{drawn with identical length at every $t$}, since that is the
content of Proposition~\ref{prop:t-invariance}.
The visual point of the figure is that the right-hand arrows do not shrink as
$t \to 1$. This replaces roughly 250 words and is the figure the reader should
remember when reaching Chapter~\ref{ch:asymmetry}.}
